% GENERATED FILE. Edit canonical lessons, then run npm run generate:books.
% canonical-source-hash: fnv1a64:99c1da568d039e01
% canonical-lessons: ZH-C412-gangkou, ZH-C412-jiaban, ZH-C412-chuancang, ZH-C412-kaoan, ZH-C412-kechuan

\chapter{Harbour and Deck}
\label{ch:zh-harbour-and-deck}
\hlchaptermodality{412}

\begin{chapteropening}
\textbf{By the end of this chapter:} \emph{I can talk about the harbour, the deck and the cabins, a boat coming in to shore and a passenger ship.}

Complete the last lesson of chapter 412: I can talk about the harbour, the deck and the cabins, a boat coming in to shore and a passenger ship.
\end{chapteropening}

\section[gǎngkǒu]{gǎngkǒu — port; harbour}

\label{lesson:ZH-C412-gangkou}



\begin{quote}
Before the new one: say the Mandarin for a ferry across the water, then the Mandarin for a life jacket.
\end{quote}

\subsection*{You'll want to know: gǎngkǒu}
\emph{gǎngkǒu} — \textquotedblleft{}port; harbour\textquotedblright{}.

\begin{quote}
\emph{zhège gǎngkǒu hěn dà} — This harbour is very big.
\end{quote}

Its characters come later; for now it is a word for the ear.

Two third tones meet in this word: gǎng is said with a rising tone before kǒu, and the spelling keeps both marks.

\begin{cousinweb}
\emph{gǎngkǒu} is built from \emph{gǎng} 'a harbour' and \zh{口} 'a mouth', both inherited from Middle Chinese, \emph{kaewngX} and \emph{khuwX} (Guangyun, in Baxter's transcription), which glosses the first as a branch of a river.
\end{cousinweb}

\begin{grammarlens}[title={What you've built}]
One more word for getting around.
\end{grammarlens}

\subsection*{Your turn}
\begin{itemize}
\raggedright
  \item \emph{Say it:} \emph{gǎngkǒu}
  \item \emph{Say it:} \emph{gǎngkǒu}, once more
  \item \emph{Say it:} \emph{zhège gǎngkǒu hěn dà} — This harbour is very big.
  \item \emph{Recall:} say the Mandarin for a ferry across the water, then the Mandarin for a life jacket, then say \emph{gǎngkǒu} again
\end{itemize}

\begin{tcolorbox}[breakable,colback=teal!4,colframe=teal!35!black,title={Before you move on}]
What does gǎngkǒu mean? (\textquotedblleft{}Port; harbour\textquotedblright{}.) Say it once more.
\end{tcolorbox}
\section[jiǎbǎn]{jiǎbǎn — ship's deck}

\label{lesson:ZH-C412-jiaban}



\begin{quote}
Before the new one: say the Mandarin for a life jacket, then the Mandarin for a harbour for ships.
\end{quote}

\subsection*{You'll want to know: jiǎbǎn}
\emph{jiǎbǎn} — \textquotedblleft{}ship's deck\textquotedblright{}.

\begin{quote}
\emph{wǒmen qù jiǎbǎn kàn dàhǎi} — Let's go up on deck and look at the sea.
\end{quote}

Its characters come later; for now it is a word for the ear.

Two third tones meet in this word: jiǎ is said with a rising tone before bǎn, and the spelling keeps both marks.

\begin{cousinweb}
\emph{jiǎbǎn} is built from \emph{jiǎ} 'a shell, armour' and \emph{bǎn} 'a board', both inherited from Middle Chinese, \emph{kaep} and \emph{paenX} (Guangyun, in Baxter's transcription), which glosses the first as armour; how the pair came to mean a ship's deck is not settled.
\end{cousinweb}

\begin{grammarlens}[title={What you've built}]
One more word for getting around.
\end{grammarlens}

\subsection*{Your turn}
\begin{itemize}
\raggedright
  \item \emph{Say it:} \emph{jiǎbǎn}
  \item \emph{Say it:} \emph{jiǎbǎn}, once more
  \item \emph{Say it:} \emph{wǒmen qù jiǎbǎn kàn dàhǎi} — Let's go up on deck and look at the sea.
  \item \emph{Recall:} say the Mandarin for a life jacket, then the Mandarin for a harbour for ships, then say \emph{jiǎbǎn} again
\end{itemize}

\begin{tcolorbox}[breakable,colback=teal!4,colframe=teal!35!black,title={Before you move on}]
What does jiǎbǎn mean? (\textquotedblleft{}Ship's deck\textquotedblright{}.) Say it once more.
\end{tcolorbox}
\section[chuáncāng]{chuáncāng — ship's cabin}

\label{lesson:ZH-C412-chuancang}



\begin{quote}
Before the new one: say the Mandarin for a harbour for ships, then the Mandarin for the ship's deck.
\end{quote}

\subsection*{You'll want to know: chuáncāng}
\emph{chuáncāng} — \textquotedblleft{}ship's cabin\textquotedblright{}.

\begin{quote}
\emph{nǐ de bāo fàng zài chuáncāng lǐ} — Leave your bag in the cabin.
\end{quote}

Its characters come later; for now it is a word for the ear.

\begin{cousinweb}
\emph{chuáncāng} is built from \zh{船} 'a boat' and \emph{cāng} 'a cabin'. The first is inherited from Middle Chinese \emph{zywen} (Guangyun, in Baxter's transcription); no dictionary we use records an older form of the second, so its history is not settled.
\end{cousinweb}

\begin{grammarlens}[title={What you've built}]
One more word for getting around.
\end{grammarlens}

\subsection*{Your turn}
\begin{itemize}
\raggedright
  \item \emph{Say it:} \emph{chuáncāng}
  \item \emph{Say it:} \emph{chuáncāng}, once more
  \item \emph{Say it:} \emph{nǐ de bāo fàng zài chuáncāng lǐ} — Leave your bag in the cabin.
  \item \emph{Recall:} say the Mandarin for a harbour for ships, then the Mandarin for the ship's deck, then say \emph{chuáncāng} again
\end{itemize}

\begin{tcolorbox}[breakable,colback=teal!4,colframe=teal!35!black,title={Before you move on}]
What does chuáncāng mean? (\textquotedblleft{}Ship's cabin\textquotedblright{}.) Say it once more.
\end{tcolorbox}
\section[kào'àn]{kào'àn — to come in to shore}

\label{lesson:ZH-C412-kaoan}



\begin{quote}
Before the new one: say the Mandarin for the ship's deck, then the Mandarin for a ship's cabin.
\end{quote}

\subsection*{You'll want to know: kào'àn}
\emph{kào'àn} — \textquotedblleft{}to come in to shore\textquotedblright{}.

\begin{quote}
\emph{chuán mǎshàng yào kào'àn le} — The boat is about to come in.
\end{quote}

Its characters come later; for now it is a word for the ear.

\begin{cousinweb}
\emph{kào'àn} is built from \emph{kào} 'to lean against' and \emph{àn} 'a shore', both inherited from Middle Chinese, \emph{khawH} and \emph{nganH} (Guangyun, in Baxter's transcription), which glosses the second as the high bank of a river; the sense 'to lean on' for the first is not given there.
\end{cousinweb}

\begin{grammarlens}[title={What you've built}]
One more word for getting around.
\end{grammarlens}

\subsection*{Your turn}
\begin{itemize}
\raggedright
  \item \emph{Say it:} \emph{kào'àn}
  \item \emph{Say it:} \emph{kào'àn}, once more
  \item \emph{Say it:} \emph{chuán mǎshàng yào kào'àn le} — The boat is about to come in.
  \item \emph{Recall:} say the Mandarin for the ship's deck, then the Mandarin for a ship's cabin, then say \emph{kào'àn} again
\end{itemize}

\begin{tcolorbox}[breakable,colback=teal!4,colframe=teal!35!black,title={Before you move on}]
What does kào'àn mean? (\textquotedblleft{}To come in to shore\textquotedblright{}.) Say it once more.
\end{tcolorbox}
\section[\texorpdfstring{\zh{客船} (kèchuán)}{客船 (kèchuán)}]{\zh{客船} (kèchuán) — passenger ship}

\label{lesson:ZH-C412-kechuan}



\begin{quote}
Before the new one: say the Mandarin for a ship's cabin, then the Mandarin for to come in to shore.
\end{quote}

\subsection*{You'll want to know: \zh{客船}}
\textbf{\zh{客船}} — \emph{kèchuán} — \textquotedblleft{}passenger ship\textquotedblright{}.

\begin{quote}
\textbf{\zh{我们坐客船去那儿}} \emph{wǒmen zuò kèchuán qù nàr} — We're going there on a passenger ship.
\end{quote}

\begin{cousinweb}
\textbf{\zh{客船}} (\emph{kèchuán}) is built from \zh{客} 'a guest' and \zh{船} 'a boat', both inherited from Middle Chinese, \emph{khaek} and \emph{zywen} (Guangyun, in Baxter's transcription).
\end{cousinweb}

\begin{grammarlens}[title={What you've built}]
One more word for getting around.
\end{grammarlens}

\subsection*{Your turn}
\begin{itemize}
\raggedright
  \item \emph{Say it:} \emph{kèchuán}
  \item \emph{Say it:} \emph{kèchuán}, once more
  \item \emph{Say it:} \emph{wǒmen zuò kèchuán qù nàr} — We're going there on a passenger ship.
  \item \emph{Recall:} say the Mandarin for a ship's cabin, then the Mandarin for to come in to shore, then say \emph{kèchuán} again
\end{itemize}

\begin{tcolorbox}[breakable,colback=teal!4,colframe=teal!35!black,title={Before you move on}]
What does \zh{客船} mean? (\textquotedblleft{}Passenger ship\textquotedblright{}.) Say it once more.
\end{tcolorbox}
